\documentclass[11pt,reqno]{amsart}
\usepackage[T1]{fontenc}
\usepackage[utf8]{inputenc}
\usepackage{lmodern}
\usepackage{amsmath,amssymb,amsfonts,amsthm,mathtools,mathrsfs}
\usepackage{enumitem}
\usepackage{microtype}
\usepackage[colorlinks=true,linkcolor=blue,citecolor=blue,urlcolor=blue]{hyperref}
\usepackage[nameinlink,capitalize]{cleveref}
\allowdisplaybreaks
\setlength{\emergencystretch}{2em}

\newtheorem{theorem}{Theorem}[section]
\newtheorem{proposition}[theorem]{Proposition}
\newtheorem{lemma}[theorem]{Lemma}
\newtheorem{corollary}[theorem]{Corollary}
\theoremstyle{definition}
\newtheorem{definition}[theorem]{Definition}
\newtheorem{assumption}[theorem]{Assumption}
\newtheorem{example}[theorem]{Example}
\theoremstyle{remark}
\newtheorem{remark}[theorem]{Remark}

\newcommand{\C}{\mathbb C}
\newcommand{\E}{\mathbb E}
\newcommand{\F}{\mathcal F}
\newcommand{\Pp}{\mathbb P}
\newcommand{\Qp}{\mathbb Q}
\newcommand{\R}{\mathbb R}
\newcommand{\T}{\mathbb T}
\newcommand{\cE}{\mathcal E}
\newcommand{\cP}{\mathcal P}
\newcommand{\Id}{\mathrm{Id}}
\newcommand{\tr}{\operatorname{tr}}
\newcommand{\esssup}{\operatorname*{ess\,sup}}
\newcommand{\norm}[1]{\left\lVert #1\right\rVert}
\newcommand{\abs}[1]{\left\lvert #1\right\rvert}

\title[Tangent-law characterization and stochastic representations]{Tangent-Law Characterization and Stochastic Representations of Theta-Semigroups}
\author{Qian Qi}
\address{Peking University, Beijing 100871, China}
\email{qiqian@pku.edu.cn}
\subjclass[2020]{Primary 60H10, 60H30; Secondary 35K55, 35D40, 49L25}
\keywords{nonlinear semigroup, tangent law, characterization theorem, Girsanov theorem, BSDE, PPDE, second-order BSDE}
\date{August 29, 2026}

\begin{document}

\begin{abstract}
We develop a representation theory that is intrinsic to the theta-semigroups
constructed from deterministic moving collision dynamics in the companion
papers.  The first result is a parabolic implicit-function theorem: under
explicit uniform parabolicity and H\"older hypotheses, the terminal-data
solution map is Fr\'echet differentiable and its derivative is the unique
solution of the linear tangent equation.  Differentiability is proved rather
than assumed.

For the actual entropic theta-semigroup, the terminal map is real analytic on
bounded payoffs.  Its first derivative is expectation under an explicit
payoff-dependent tangent law and its second derivative is the corresponding
covariance.  The tangent kernels form a cocycle obtained by differentiating
nonlinear time consistency.  We prove a converse characterization: a monotone
cash-additive $C^2$ functional whose tangent probabilities have the covariance
curvature
\[
 D^2\mathfrak E_\phi(f,g)
 =\vartheta\operatorname{Cov}_{Q_\phi}(f,g)
\]
must be the logarithmic moment-generating functional of its zero-payoff tangent
law.  The dynamic version characterizes entropic theta-semigroups by their
tangent-kernel bundle.

The finite collision recursions of the homogenization paper have exact
pathwise exponential tangent laws.  We prove convergence of these microscopic
tangent laws, and of their first two derivative operators, to the continuum
tangent diffusion.  For smooth data the density process is a stochastic
exponential with integrand $\vartheta\sigma^TDu$.  A bounded-martingale/BMO
argument and, in the bounded-gradient window, Novikov's condition give a
genuine Girsanov theorem with a Radon--Nikodym density and changed Brownian
drift.  No inversion of $Z=\sigma^TDu$ is used.

The same solution admits a quadratic BSDE, a linear tangent BSDE under the
tilted law, a path-dependent PPDE/BSDE extension, an exact pure energy-game
representation, and a stable volatility-control 2BSDE branch.
\end{abstract}

\maketitle
\tableofcontents

\section{Introduction}

A genuinely nonlinear semigroup cannot be represented by one
payoff-independent Markov kernel for all terminal payoffs.  It may nevertheless
carry a canonical bundle of payoff-dependent tangent laws.  For the entropic
theta-semigroup this bundle has three rigid features: the first derivative is a
probability kernel, the second derivative is its covariance multiplied by one
constant $\vartheta$, and the kernels compose according to nonlinear time
consistency.  We prove that these properties characterize the semigroup.

This paper begins after the deterministic response, rough homogenization, and
HJB limits have been established.  It does not use stochastic representation
to derive those limits.  Instead, it proves that the finite deterministic
collision evaluations and their derivatives converge to the continuum
stochastic objects.

Classical BSDE theory begins with Pardoux--Peng \cite{PardouxPeng1990};
quadratic equations are treated in \cite{Kobylanski2000}.  Functional It\^o
calculus and PPDEs are developed in
\cite{Dupire2009,ContFournie2013,EkrenTouziZhang2014I,
EkrenTouziZhang2014II}.  Nondominated second-order BSDEs are treated in
\cite{CheriditoSonerTouziVictoir2007,SonerTouziZhang2012}.  Our contribution is
the tangent-law characterization, its microscopic convergence from one
collision system, and the dynamically consistent Girsanov/BSDE bundle.

\section{The fixed-law obstruction}

Let $E$ be a Polish state space.

\begin{theorem}[A payoff-independent law forces linearity]\label{thm:fixed-law}
Suppose that for fixed $s<t$ and every initial state $x$ there is one
probability kernel $P_{s,t}(x,dy)$ such that, for every
$\phi\in C_b(E)$,
\[
 \cE_{s,t}[\phi](x)=\int_E\phi(y)P_{s,t}(x,dy).
\]
Then $\cE_{s,t}$ is linear, positive, and constant preserving.  Hence a
semigroup with a genuinely nonlinear generator cannot have one universal
payoff-independent classical law.
\end{theorem}

\begin{proof}
Integration against a fixed probability kernel preserves arbitrary linear
combinations, order, and constants at every initial state.
\end{proof}

The theorem does not exclude payoff-dependent tangent laws, controlled
families, games, or nondominated representations.

\section{A parabolic implicit-function theorem}

Work on the torus $\T^d$.  Fix $0<\alpha<1$ and put
\[
 \mathbb X=C^{1+\alpha/2,2+\alpha}([0,T]\times\T^d),
 \qquad
 \mathbb Y=C^{\alpha/2,\alpha}([0,T]\times\T^d)
 \times C^{2+\alpha}(\T^d).
\]
Let $F(t,x,r,p,X)$ be $C^k$, $k\ge2$, on a neighbourhood of the relevant jet
set and consider
\begin{equation}\label{eq:PDE}
 -u_t-F(t,x,u,Du,D^2u)=0,
 \qquad u(T)=\phi.
\end{equation}

\begin{assumption}[Uniform classical window]\label{ass:window}
There are $0<\lambda\le\Lambda$ such that
\[
 \lambda I\le F_X(t,x,r,p,X)\le\Lambda I
\]
on the jet neighbourhood.  All derivatives of $F$ through order $k$ are
bounded and H\"older there.  At one pair $(u_0,\phi_0)$,
$u_0\in\mathbb X$ solves \eqref{eq:PDE}.
\end{assumption}

\begin{theorem}[Fr\'echet differentiability of the terminal map]\label{thm:IFT}
Under \cref{ass:window}, there are neighbourhoods of $\phi_0$ in
$C^{2+\alpha}$ and of $u_0$ in $\mathbb X$ on which the solution map
$\phi\mapsto u[\phi]$ is $C^{k-1}$.  For a terminal direction $\eta$, the first
derivative $v=Du[\phi]\eta$ is the unique classical solution of
\begin{equation}\label{eq:tangent-general}
 -v_t-F_rv-F_p\cdot Dv-F_X:D^2v=0,
 \qquad v(T)=\eta,
\end{equation}
where the coefficients are evaluated along $u[\phi]$.
\end{theorem}

\begin{proof}
Define
\[
 \mathfrak F(u,\phi)=
 \bigl(-u_t-F(t,x,u,Du,D^2u),\ u(T)-\phi\bigr).
\]
H\"older composition makes $\mathfrak F:\mathbb X\times C^{2+\alpha}	o
\mathbb Y$ a $C^{k-1}$ map.  Its derivative in $u$ is the uniformly parabolic
terminal operator
\[
 v\mapsto
 \bigl(-v_t-F_rv-F_p\cdot Dv-F_X:D^2v,\ v(T)\bigr).
\]
Schauder theory makes this a bounded isomorphism from $\mathbb X$ to
$\mathbb Y$.  The Banach implicit-function theorem gives the $C^{k-1}$
solution map.  Differentiating
$\mathfrak F(u[\phi],\phi)=0$ yields \eqref{eq:tangent-general}.
\end{proof}

No conclusion about a degenerate regularization limit is inferred from local
uniform convergence alone.  Whenever a degenerate tangent limit is used, its
uniform tangent estimates are stated separately.

\section{The actual entropic theta-semigroup}

Let $X^{t,x}$ solve
\begin{equation}\label{eq:forward}
 dX_s=b(s,X_s)ds+\sigma(s,X_s)dW_s,
 \qquad X_t=x,
\end{equation}
where $b,\sigma$ are bounded globally Lipschitz and
$a=\sigma\sigma^T$ is uniformly elliptic.  Let $P_{t,T}$ be the transition
operator.  For $\vartheta\ne0$ define
\begin{equation}\label{eq:theta}
 \cE_{t,T}^{\vartheta}[\phi](x)
 ={1\over\vartheta}\log P_{t,T}(e^{\vartheta\phi})(x).
\end{equation}

\begin{theorem}[Global analytic terminal map]\label{thm:analytic}
For fixed $(t,x)$, the map
$C_b(\R^d)\ni\phi\mapsto\cE_{t,T}^{\vartheta}[\phi](x)$ is real analytic.
Its first two Fr\'echet derivatives are
\begin{align}
 D\cE_{t,T}^{\vartheta}[\phi]\eta
 &=\frac{P_{t,T}(e^{\vartheta\phi}\eta)}
        {P_{t,T}(e^{\vartheta\phi})},\label{eq:first}\\
 D^2\cE_{t,T}^{\vartheta}[\phi](\eta_1,\eta_2)
 &=\vartheta\left[
 \frac{P_{t,T}(e^{\vartheta\phi}\eta_1\eta_2)}
      {P_{t,T}(e^{\vartheta\phi})}
 -\prod_{j=1}^2
 \frac{P_{t,T}(e^{\vartheta\phi}\eta_j)}
      {P_{t,T}(e^{\vartheta\phi})}
 \right].\label{eq:second}
\end{align}
Higher derivatives are the tilted joint cumulants multiplied by
$\vartheta^{n-1}$.
\end{theorem}

\begin{proof}
On every bounded ball of $C_b$, the exponential Nemytskii map has an
absolutely convergent power series in the sup norm.  The Markov operator is
bounded and positive, and
\[
 P_{t,T}e^{\vartheta\phi}
 \ge e^{-\abs\vartheta\norm\phi_\infty}>0.
\]
The logarithm is analytic on a neighbourhood of the range.  Differentiating
the composition gives \eqref{eq:first}--\eqref{eq:second}; higher logarithmic
derivatives are cumulants.
\end{proof}

\section{Tangent laws and their cocycle}

Define the terminally tilted law
\begin{equation}\label{eq:tilted-law}
 {d\Qp_{t,x}^{\phi}\over d\Pp_{t,x}}
 ={e^{\vartheta\phi(X_T)}
 \over P_{t,T}e^{\vartheta\phi}(x)}.
\end{equation}
Then \eqref{eq:first} is
\begin{equation}\label{eq:tangent-expectation}
 D\cE_{t,T}^{\vartheta}[\phi]\eta(x)
 =\E^{\Qp_{t,x}^{\phi}}\eta(X_T).
\end{equation}
For an intermediate interval define
\begin{equation}\label{eq:tangent-kernel}
 K_{s,t}^{\psi}(x,dy)=
 {e^{\vartheta\psi(y)}P_{s,t}(x,dy)
 \over P_{s,t}e^{\vartheta\psi}(x)}.
\end{equation}

\begin{theorem}[Tangent-kernel cocycle]\label{thm:cocycle}
The nonlinear semigroup is time consistent and its tangent kernels satisfy
\begin{equation}\label{eq:cocycle}
 K_{r,t}^{\phi}
 =K_{r,s}^{\cE_{s,t}^{\vartheta}[\phi]}
  K_{s,t}^{\phi},
 \qquad r\le s\le t.
\end{equation}
Equivalently,
\[
 D\cE_{r,t}^{\vartheta}[\phi]
 =D\cE_{r,s}^{\vartheta}
   [\cE_{s,t}^{\vartheta}[\phi]]
 \circ D\cE_{s,t}^{\vartheta}[\phi].
\]
\end{theorem}

\begin{proof}
The Markov property gives
\[
 P_{r,t}e^{\vartheta\phi}
 =P_{r,s}\exp\{\vartheta\cE_{s,t}^{\vartheta}[\phi]\}.
\]
This is nonlinear time consistency.  The Radon--Nikodym factors in
\eqref{eq:tangent-kernel} multiply and the intermediate normalizations cancel,
giving \eqref{eq:cocycle}.  The derivative identity follows from the Banach
chain rule.
\end{proof}

\section{A characterization by tangent curvature}

Let $(E,\mathcal E)$ be a measurable space and let
$\mathfrak E:L^\infty(E)\to\R$ be twice Fr\'echet differentiable.

\begin{assumption}[Probability tangents and covariance curvature]\label{ass:curvature}
The functional is monotone, cash additive, and $\mathfrak E(0)=0$.  For every
$\phi$, there is a probability measure $Q_\phi$ such that
\begin{align}
 D\mathfrak E(\phi)f&=Q_\phi(f),\label{eq:tangent-prob}\\
 D^2\mathfrak E(\phi)(f,g)
 &=\vartheta\left(Q_\phi(fg)-Q_\phi(f)Q_\phi(g)\right)
 \label{eq:cov-curvature}
\end{align}
for one fixed nonzero constant $\vartheta$.
\end{assumption}

\begin{theorem}[Static tangent-law characterization]\label{thm:characterization}
Under \cref{ass:curvature},
\begin{equation}\label{eq:characterization}
 \mathfrak E(\phi)
 ={1\over\vartheta}\log Q_0(e^{\vartheta\phi}),
 \qquad
 {dQ_\phi\over dQ_0}
 ={e^{\vartheta\phi}\over Q_0(e^{\vartheta\phi})}.
\end{equation}
Thus the tangent-probability and covariance-curvature identities characterize
the logarithmic moment-generating functional; they are not merely consequences
of a chosen name.
\end{theorem}

\begin{proof}
Fix $\phi$ and put $Q_t=Q_{t\phi}$.  For every bounded $f$,
\begin{align*}
 {d\over dt}Q_t(f)
 &=D^2\mathfrak E(t\phi)(\phi,f)\\
 &=\vartheta\left(Q_t(\phi f)-Q_t(\phi)Q_t(f)\right).
\end{align*}
The probability measures
\[
 \widetilde Q_t(f)=
 {Q_0(e^{\vartheta t\phi}f)
  \over Q_0(e^{\vartheta t\phi})}
\]
solve the same weak differential equation and have the same initial condition.
Uniqueness follows by Gronwall after testing the difference against bounded
functions.  Hence $Q_t=\widetilde Q_t$.  Finally,
\[
 {d\over dt}\mathfrak E(t\phi)=Q_t(\phi)
 ={d\over dt}{1\over\vartheta}
 \log Q_0(e^{\vartheta t\phi}).
\]
Integrate from zero to one.
\end{proof}

\begin{theorem}[Dynamic characterization]\label{thm:dynamic-characterization}
Let $(\mathfrak E_{s,t})_{s\le t}$ be a monotone cash-additive $C^2$ dynamic
semigroup.  Suppose that for every $(s,t,x)$ its first and second terminal
derivatives satisfy \eqref{eq:tangent-prob}--\eqref{eq:cov-curvature}, that
the zero-payoff tangent kernels $P_{s,t}$ form a Markov semigroup, and that the
tangent kernels satisfy the cocycle \eqref{eq:cocycle}.  Then
\[
 \mathfrak E_{s,t}[\phi](x)
 ={1\over\vartheta}\log P_{s,t}e^{\vartheta\phi}(x).
\]
\end{theorem}

\begin{proof}
Apply \cref{thm:characterization} at each $(s,t,x)$.  The base probability is
the zero-payoff tangent $P_{s,t}(x,dy)$.  Its Markov property and the tangent
cocycle give dynamic time consistency.
\end{proof}

\section{Microscopic tangent laws converge}

Let $X^h$ be the finite collision recursion from the homogenization paper
\cite{QiRoughRound4}, with path law $P_h^{t,x}$ obtained by pushing the
Lebesgue fast reference law through the deterministic collision map.  Define
\begin{equation}\label{eq:discrete-theta}
 \cE_{h,t,T}^{\vartheta}[\phi](x)
 ={1\over\vartheta}\log
 \E_{P_h^{t,x}}e^{\vartheta\phi(X_T^h)}
\end{equation}
and
\begin{equation}\label{eq:discrete-tangent}
 {dQ_{h,t,x}^{\phi}\over dP_h^{t,x}}
 ={e^{\vartheta\phi(X_T^h)}
  \over\E_{P_h^{t,x}}e^{\vartheta\phi(X_T^h)}}.
\end{equation}

\begin{theorem}[Microscopic-to-continuum tangent convergence]\label{thm:micro-tangent}
Assume $X^h\Rightarrow X$ in path space, as proved for the moving collision
system in the companion paper, and let $\phi$ be bounded continuous.  Then
\begin{enumerate}[label=(\roman*)]
\item $Q_{h,t,x}^{\phi}\Rightarrow Q_{t,x}^{\phi}$, where
\[
 {dQ_{t,x}^{\phi}\over dP_{t,x}}
 ={e^{\vartheta\phi(X_T)}
  \over\E_{P_{t,x}}e^{\vartheta\phi(X_T)}};
\]
\item for every bounded continuous terminal direction $\eta$,
\[
 D\cE_{h,t,T}^{\vartheta}[\phi]\eta
 \longrightarrow D\cE_{t,T}^{\vartheta}[\phi]\eta;
\]
\item the second derivatives converge to the continuum tilted covariances;
\item the discrete tangent kernels converge jointly on every finite time
partition and their cocycle converges to \eqref{eq:cocycle}.
\end{enumerate}
\end{theorem}

\begin{proof}
For every bounded continuous path functional $G$,
\[
 \E_{Q_h^\phi}G
 ={\E_{P_h}[G e^{\vartheta\phi(X_T^h)}]
  \over\E_{P_h}e^{\vartheta\phi(X_T^h)}}.
\]
Weak convergence and bounded continuity give convergence of numerator and
denominator, and the denominator is uniformly positive.  This proves (i).
The first and second derivative formulas are expectations and covariances under
$Q_h^\phi$, giving (ii) and (iii).  Apply the same argument to joint path
functionals at finitely many times; the exact discrete tower identity passes to
the limit, proving (iv).
\end{proof}

This theorem is the direct deterministic-microscopic origin of the tangent
law.  It is not a post hoc representation of an unrelated diffusion.

\section{The tangent PDE}

Assume now that $b,\sigma$, and $\phi$ are smooth enough for a classical
solution and set
$u(t,x)=\cE_{t,T}^{\vartheta}[\phi](x)$.  Then
\begin{equation}\label{eq:entropic-PDE}
 -u_t-b\cdot Du-{1\over2}a:D^2u
 -{\vartheta\over2}\abs{\sigma^TDu}^2=0,
 \qquad u(T)=\phi.
\end{equation}
For a terminal perturbation $\eta$, put
$v=D\cE_{t,T}^{\vartheta}[\phi]\eta$.

\begin{theorem}[Tangent PDE]\label{thm:tangent-PDE}
The derivative $v$ is the unique solution of
\begin{equation}\label{eq:tangent-PDE}
 -v_t-\left(b+\vartheta aDu\right)\cdot Dv
 -{1\over2}a:D^2v=0,
 \qquad v(T)=\eta.
\end{equation}
Its transition law is precisely the tangent kernel
$K_{s,t}^{\cE_{t,T}^{\vartheta}[\phi]}$.
\end{theorem}

\begin{proof}
Apply \cref{thm:IFT} to \eqref{eq:entropic-PDE}; differentiating the quadratic
gradient term gives $\vartheta aDu\cdot Dv$.  The Feynman--Kac formula for the
linear equation identifies its transition kernel.  Formula
\eqref{eq:tangent-expectation} identifies the same kernel as the tilted
tangent law.
\end{proof}

\section{A genuine Girsanov theorem}

Set
\[
 Z_s=\sigma(s,X_s)^TDu(s,X_s),
 \qquad
 M_s={e^{\vartheta u(s,X_s)}
 \over e^{\vartheta u(t,x)}}.
\]
The linearized exponential $e^{\vartheta u}$ solves the backward Kolmogorov
equation.

\begin{theorem}[Tangent-law Girsanov theorem]\label{thm:Girsanov}
Let $\phi$ be bounded.  Then $M$ is a positive uniformly integrable martingale,
$M_t=1$, and
\begin{equation}\label{eq:stochastic-exponential}
 M_s=\mathcal E\left(
 \int_t^s\vartheta Z_r\,dW_r\right).
\end{equation}
It defines the tangent law on $\F_s$ by
$d\Qp^\phi/d\Pp=M_s$.  Under $\Qp^\phi$,
\begin{equation}\label{eq:Q-Brownian}
 W_s^{\Qp}=W_s-\int_t^s\vartheta Z_rdr
\end{equation}
is Brownian, and the state has drift
$b+\vartheta aDu$.  If $Z$ is bounded, Novikov's condition
\[
 \E\exp\left({\vartheta^2\over2}\int_t^T\abs{Z_s}^2ds\right)<\infty
\]
holds directly.  For bounded terminal data the bounded density martingale also
gives the BMO/uniform-integrability route when a global gradient bound is not
used.
\end{theorem}

\begin{proof}
It\^o's formula and \eqref{eq:entropic-PDE} give
$dM_s=M_s\vartheta Z_s\,dW_s$.  Since
\[
 M_s={\E[e^{\vartheta\phi(X_T)}\mid\F_s]
 \over\E e^{\vartheta\phi(X_T)}},
\]
it is positive, bounded above and below by constants depending only on
$\norm\phi_\infty$, and uniformly integrable.  Therefore
\eqref{eq:stochastic-exponential} defines a probability change of measure.
Girsanov's theorem yields \eqref{eq:Q-Brownian}; substitution in
\eqref{eq:forward} gives the changed drift.  The bounded-$Z$ statement is
Novikov's criterion, while the bounded positive density gives the alternative
BMO route.
\end{proof}

This is a measure-change theorem with a Radon--Nikodym derivative; it is not an
algebraic shift of the Hamiltonian.

\section{Quadratic and tangent BSDEs}

\begin{theorem}[Quadratic BSDE representation]\label{thm:BSDE}
Let $Y_s=u(s,X_s)$ and $Z_s=\sigma^TDu(s,X_s)$.  Then
\begin{equation}\label{eq:BSDE}
 Y_s=\phi(X_T)+
 \int_s^T{\vartheta\over2}\abs{Z_r}^2dr
 -\int_s^TZ_r\,dW_r.
\end{equation}
For bounded terminal data this quadratic BSDE has the entropic solution
\eqref{eq:theta}.
\end{theorem}

\begin{proof}
Apply It\^o's formula to $u(s,X_s)$ and use
\eqref{eq:entropic-PDE}.  Exponentiating $\vartheta Y$ gives the bounded
martingale $P_{s,T}e^{\vartheta\phi}$, which identifies the solution.
\end{proof}

\begin{corollary}[Linear tangent BSDE]\label{cor:tangent-BSDE}
For a terminal direction $\eta$, the tangent process
$V_s=v(s,X_s)$ satisfies under $\Qp^\phi$
\begin{equation}\label{eq:tangent-BSDE}
 V_s=\eta(X_T)-\int_s^T\widehat Z_r\,dW_r^{\Qp}.
\end{equation}
Thus the derivative semigroup is the martingale representation under the
payoff-dependent tangent law.
\end{corollary}

\begin{proof}
The drift in the tangent PDE \eqref{eq:tangent-PDE} is exactly the drift of
$X$ under $\Qp^\phi$.  It\^o's formula gives the martingale identity.
\end{proof}

No inversion of the nonlinear map $p\mapsto\sigma(x,p)^Tp$ appears: in this
actual branch $\sigma$ is a state coefficient, and the tangent drift is written
directly in terms of the already proved decoupling field.

\section{Path-dependent entropic evaluation}

Let $\Omega=C([0,T];\R^d)$ and let $\Phi:\Omega\to\R$ be bounded uniformly
continuous.  Define
\begin{equation}\label{eq:path-evaluation}
 U(t,\omega)={1\over\vartheta}
 \log\E_{t,\omega}e^{\vartheta\Phi(X)}.
\end{equation}

\begin{theorem}[Path DPP, PPDE, and BSDE]\label{thm:path}
The path evaluation is time consistent.  In the smooth functional It\^o
window it solves
\begin{equation}\label{eq:PPDE}
 -\partial_tU-b\cdot\partial_\omega U
 -{1\over2}a:\partial_{\omega\omega}^2U
 -{\vartheta\over2}
 \abs{\sigma^T\partial_\omega U}^2=0,
 \qquad U(T,\omega)=\Phi(\omega).
\end{equation}
It has the quadratic BSDE representation
\[
 Y_s=\Phi(X)+\int_s^T{\vartheta\over2}\abs{Z_r}^2dr
 -\int_s^TZ_r\,dW_r.
\]
For bounded uniformly continuous data, \eqref{eq:path-evaluation} is the unique
bounded path-viscosity solution in the standard comparison class.
\end{theorem}

\begin{proof}
The tower property of conditional exponential moments gives the path DPP.
Functional It\^o calculus gives \eqref{eq:PPDE} and the BSDE in the smooth
window.  The exponential transform turns the PPDE into the linear path
Kolmogorov equation.  Stability and comparison for bounded path-viscosity
solutions identify the general case; see
\cite{ContFournie2013,EkrenTouziZhang2014I,EkrenTouziZhang2014II}.
\end{proof}

\section{The pure energy path game}

Consider
\[
 dX_s=(u_s+v_s)ds+\sigma dW_s
\]
with running payoff
$-\mu u_s^2/2+\nu v_s^2/2$ and bounded path terminal payoff $\Phi(X)$.
Set
\[
 c={1\over2}(\mu^{-1}-\nu^{-1}),
 \qquad \vartheta={2c\over\sigma^2},
\]
and assume $c\ne0$.

\begin{theorem}[Exact pure path-game representation]\label{thm:path-game}
The lower and upper path-game values coincide and equal
\begin{equation}\label{eq:path-game-value}
 U(t,\omega)={1\over\vartheta}
 \log\E_{t,\omega}e^{\vartheta\Phi(X)}.
\end{equation}
In the smooth window the unique pure feedback saddle is
\[
 u_s^*={\partial_\omega U(s,X)\over\mu},
 \qquad
 v_s^*=-{\partial_\omega U(s,X)\over\nu}.
\]
The same value is the bounded path-viscosity solution and the quadratic BSDE
solution of \cref{thm:path}.
\end{theorem}

\begin{proof}
The local Hamiltonian is $c p^2$.  With $p=\partial_\omega U$,
\[
 -cp^2+p(u+v)-{\mu\over2}u^2+{\nu\over2}v^2
 =-{\mu\over2}(u-p/\mu)^2
 +{\nu\over2}(v+p/\nu)^2.
\]
Functional It\^o calculus therefore gives
$J(u,v^*)\le U\le J(u^*,v)$ for every admissible pair.  The Cole--Hopf
transform with $\vartheta=2c/\sigma^2$ gives
\eqref{eq:path-game-value}.  Approximation and path-viscosity comparison pass
the identity beyond the smooth window.
\end{proof}

\section{A nondominated second-order branch}

Let $0<\underline\sigma\le\overline\sigma$ and let
$\mathcal A$ be the progressively measurable volatility controls with values
in this interval.  Put
\[
 X_s^{\alpha}=x+\int_t^s\alpha_r\,dW_r
\]
and define
\begin{equation}\label{eq:G-evaluation}
 \mathcal E_{t,T}^G[\Phi](\omega)
 =\sup_{\alpha\in\mathcal A}
 \E\Phi\left(\omega\otimes_tX^\alpha\right).
\end{equation}
The family of induced laws is stable under conditioning and pasting.

\begin{theorem}[Path $G$-heat equation and 2BSDE]\label{thm:2BSDE}
The evaluation \eqref{eq:G-evaluation} is time consistent and is the unique
bounded path-viscosity solution of
\begin{equation}\label{eq:G-PPDE}
 -\partial_tU-G(\partial_{\omega\omega}^2U)=0,
 \qquad
 G(\Gamma)={1\over2}
 \sup_{a\in[\underline\sigma^2,\overline\sigma^2]}a\Gamma.
\end{equation}
It admits the nondominated second-order BSDE representation
\begin{equation}\label{eq:2BSDE}
 Y_s=\Phi-\int_s^TZ_r\,dB_r+K_T^P-K_s^P,
 \qquad P\text{-a.s. for every }P\in\mathcal P,
\end{equation}
with the usual minimality condition on the increasing family $K^P$.
\end{theorem}

\begin{proof}
Stability under conditioning and pasting gives the DPP.  Dynamic programming
and the functional It\^o expansion give \eqref{eq:G-PPDE}.  The representation
by the essential supremum of continuation laws and the minimal increasing
process is the standard 2BSDE construction
\cite{CheriditoSonerTouziVictoir2007,SonerTouziZhang2012}.  Comparison
identifies the value.
\end{proof}

\section{Conclusion}

The representation layer now contains a theorem that is specific to the theta
structure: tangent probabilities with covariance curvature characterize the
entropic functional, and their dynamic cocycle characterizes the semigroup.
These tangent laws are limits of exact exponential tilts of the finite
deterministic collision recursions.  The Girsanov theorem has an actual density
and changed Brownian motion, the terminal derivative is proved by a parabolic
implicit-function theorem, and the path and second-order branches are typed by
their own DPPs.

\section*{Disclosure and review status}
The author is responsible for all proofs, references, and computations.
Structural verifiers do not certify the analytical arguments.  This revision
is prepared for a new independent formal review.

\bibliographystyle{plain}
\bibliography{references-round4}

\end{document}
